\documentclass[12pt]{article}
\usepackage[english]{babel}
\usepackage{natbib}
\usepackage{url}
\usepackage[utf8x]{inputenc}
\usepackage{amsmath}
\usepackage{graphicx}
\graphicspath{{images/}}
\usepackage{parskip}
\usepackage{fancyhdr}
\usepackage{vmargin}
\usepackage{longtable}
\usepackage{multirow}

%\usepackage{opensans}% 
\renewcommand{\familydefault}{\sfdefault}

\usepackage[normalem]{ulem}
\setmarginsrb{3 cm}{2.5 cm}{3 cm}{2.5 cm}{1 cm}{1.5 cm}{1 cm}{1.5 cm}

\title{Task 2 Presubmission}								% Title
\author{Group 5}								% Author
\date{\today}											% Date

\makeatletter
\let\thetitle\@title
\let\theauthor\@author
\let\thedate\@date
\makeatother

\pagestyle{fancy}
\fancyhf{}
\rhead{\theauthor}
\lhead{\thetitle}
\cfoot{\thepage}

\begin{document}


%%%%%%%%%%%%%%%%%%%%%%%%%%%%%%%%%%%%%%%%%%%%%%%%%%%%%%%%%%%%%%%%%%%%%%%%%%%%%%%%%%%%%%%%%

\begin{titlepage}
	\centering
    \vspace*{0.5 cm}
    \includegraphics[scale = 0.75]{UCT.jpg}\\[1.0 cm]	% University Logo
    \LARGE University of Cape Town\\[1.0 cm]	% University Name
	\Large EEE3088F\\				% Course Code
	\large Group 5\\				% Course Name
	\rule{\linewidth}{0.2 mm} \\[0.4 cm]
	{ \huge \bfseries \thetitle}\\
	\rule{\linewidth}{0.2 mm} \\[0.5 cm]
	
	\begin{minipage}{0.4\textwidth}
		\begin{flushleft} \large
			\emph{Authors:}\\
      Torsten Babl (leader)\\
      Ebrahim Allie\\
      Yong Hao Chen\\
      Yashiv Fakir\\
      Adam Gild\\
      Jason Hillebrand\\
      Lindokuhle Lubisi\\
      Iviwe Malotana\\
      Samuel Mbiya
      
			\end{flushleft}
			\end{minipage}~
			\begin{minipage}{0.4\textwidth}
			\begin{flushright} \large
			\emph{Student Number:} \\
      BBLTOR001\\
      ALLEBR004\\
      CHNYON001\\
      FKRYAS002\\
      GLDADA002\\
      HLLJAS007\\
      LBSLIN008\\
      MLTIVI001\\
      MBYSAM003\\									% Your Student Number
		\end{flushright}
	\end{minipage}\\[0.5 cm]
	
	{\large \thedate}\\[0.5 cm]
 
	\vfill
	
\end{titlepage}

%%%%%%%%%%%%%%%%%%%%%%%%%%%%%%%%%%%%%%%%%%%%%%%%%%%%%%%%%%%%%%%%%%%%%%%%%%%%%%%%%%%%%%%%%

\tableofcontents
\pagebreak

%%%%%%%%%%%%%%%%%%%%%%%%%%%%%%%%%%%%%%%%%%%%%%%%%%%%%%%%%%%%%%%%%%%%%%%%%%%%%%%%%%%%%%%%%

\section{Executive Summary}

\section{Group-member Contributions}
\newpage
\section{User Requirements and Acceptance Test Procedures}

The aim of the commissioned sensing and reporting system is to identify and report faults on 33kV transmission lines 
that do not trigger an earth leakage or overcurrent fault condition on existing switching systems, yet still pose 
a safety risk to the public. Such faults are often caused by storm damage, motor/construction vehicle impact and 
malicious damage to public property generally caused by theft. In addition to post-fault reporting the early recognition
of structural degradation can potentially save the Utility money by allowing planned and targeted maintenance to prevent 
future unexpected faults.

Below follows the Utility's user requirements and the acceptance tests thereof.

\subsection{User stories}
The commissioned product is envisioned to fulfil the following tasks:\newline
\emph{See the following subsections for detailed scope of key terms.}
\begin{itemize}
  \item The product must report a fault occurring along the transmission line shortly after it occurs.
  \item The product must, upon request, report data useful to assess the structural integrity of the pole.
\end{itemize}

\subsection{User Requirements of The Utility}
The following faults must be detectable:

\begin{center}
  \begin{table}[htp!]
    \caption{Table of Faults to Detect}
    
    \hskip-2.2cm\begin{tabular}{|p{2cm}|p{4cm}|p{8cm}|p{4cm}|}
        \hline
        \textbf{Fault No.} & \textbf{Fault Name} & \textbf{Detailed Description} & \textbf{Illustration} \\
        \hline
        F01 & Phase failure & Loss  of  energy  of  one  or  more  lines  due  to  faults
        between the node and the supply. & - \\\hline

        F02 & Line detachment & One or more lines are detached from the isolator on the A-frame while the pole remains
        upright. & - \\\hline

        F03 & Pole dislodgment and/or rotation & Pole is no longer upright or loses orientation through forces on the
        lines or pole itself. & - \\\hline

        F00 & Strong Impact to pole & An unusually strong impact to the pole is observed. & - \\\hline

        F00 & Sudden movement & The top of the pole is still upright and oriented as before but has moved out of
        position. i.e. Top segment of pole is suspended by the lines. & - \\\hline
        %Don't like these two
        F04 & Possible structural fatigue of pole & The structure of the pole is compromised due to material fatigue. & - \\\hline

        F06 & Pole and node failure due to damaging force & A damaging force is strong enough to damage both the pole
        and the node. & - \\\hline
  
    \end{tabular}    
  
  \label{tab:faults}
  \end{table}
\end{center}

\begin{center}
  \begin{table}[htp!]
    \caption{Table of General User Requirements}
    
    \hskip-2.2cm\begin{tabular}{|p{3cm}|p{2cm}|p{4cm}|p{9cm}|}
        \hline
        \textbf{Requirement Category} & \textbf{Req. No.} & \textbf{Requirement Name} & \textbf{Detailed Description} \\
        \hline
        One & R00 & Reporting time & A detected fault must be reported within 10 minutes of the incident \\\hline
        One & R00 & Product lifespan & 90\% of deployed devices must still be functional after 10 years of installation \\\hline
        %Add BOM to list of definitions
        One & R00 & Device cost & The BOM cost of a single device should not exceed 250 ZAR \\\hline
        One & R00 & Additional infrastructure & If any additional infrastructure is needed it should not add significantly
        to the cost of rollout \\\hline
        %Add ISM band to list of Definitions
        One & R00 & Communication bands & All wireless communications should use ISM bands permitted in South Africa
        to ease certification and decrease cost. \\\hline
  
    \end{tabular}    
  
  \label{tab:faults}
  \end{table}
\end{center}

\subsection{Acceptance Test procedures for User Requirements}
The user requirements set above shall be verified according to the following test criteria.


\section{System Design and Overview}

\section{Sensing Subsystem Design}

\section{Microcontroller Subsystem Design}

\section{Wireless Communications Design}

\subsection{Wireless Protocol}

The Weightless wireless technology standard has been selected for the wireless communication needs of each sensor node. Weightless supports data transmission over long distances and is a standard designed for low-power wide area networks (LPWAN) involving base stations and end devices. For the purposes of the project, some criteria was identified to measure the performance of the wireless technology standard selected: the power efficiency, range, network capacity and  functionality and ease of configurability.(complexity) Weightless-P is a open standard and operates across the range of ISM, sub GHz  frequency bands for ease of development of IoT systems. 

It was recognised that interference and noise in urban areas compared to rural areas differ, therefore the transmission range would also be affected in different ways. That's why a wireless standard with an adaptive data rate and transmission power was selected to cater for the different signal paths and the environments in which the base stations and sensors nodes will be placed. That together with the fact that the end devices using Weightless can also be configured to be idle/turn to sleep mode, will allow for better power efficiency and aid to maximise the lifespan of our sensor nodes. 
To also save on costs we also identified the need for a standard that has large network capacity, allowing the use of fewer base stations that can each cover communication between multiple sensor nodes.

[https://www.ubiik.com/post/2016/09/21/what-is-weightless-p], [https://www.electronicdesign
.com/industrial-automation/article/21807028/choosing-the-right-wireless-protocol-for-your-iot-application], [https://radiobridge.com/blog/choosing-the-right-wireless-standard-for-your-iot-sensors]

\subsection{Message Packet Design}

\subsection{Range and Error rate simulations}
[https://www.mathworks.com/help/comm/examples/bluetooth-low-energy-bit-error-rate-simulation.html], %Bit Error Rates - Need to find one for  Weightless%
[https://www.mathworks.com/help/comm/examples.html?category=rf-propagation]%Range and Power Simulation %

\section{Power Supply Design}
\subsection{User requirements in regards to power supply of system}
The device needs to function 24/7 throughout the year. South Africa is a country with several issues about their power distribution, therefore load shedding would be a major concern when the device is running. This taken in we would use a solar panel and it would charge a battery during daytime. it will have enough energy to supply the device during night time or any other load shedding issues.

The device needs to be easy to maintain throughout the service years and easy to install without requiring too much space. The user would be able to swap batteries and maintain the power supply system for a long period of time.

Since the power supply will be independent from  excess power supplies to maintain the requirement of providing power despite any disruption.

\subsection{Solar energy reliability}
Since the design chose to use a solar-energy system to charge and power the device, for the device to function at its most sufficient power, it relies heavily on the amount of sun light it receives. The main reason the design chooses a solar energy power supply system is because South Africa is one of the best places to use this technology due to its sunny climate. "As Professor Detlev Kröger, Senior Researcher and Emeritus Professor in the Department of Mechanical and Mechatronic Engineering at Stellenbosch University claims that South Africa is one of the best places in the world to develop solar power. "Citing
https://sola.africa/news/south-africas-solar-resource-compared-to-the-rest-of-the-world/

According to the government website on energy, it states that "Most areas in South Africa average more than 2 500 hours of sunshine per year, and average solar-radiation levels range between 4.5 and 6.5kWh/m2 in one day. " This basically means that the whole country is has sun light most of the time through out the year. "The annual 24-hour global solar radiation average is about 220 W/m2 for South Africa, compared with about 150 W/m2 for parts of the USA, and about 100 W/m2 for Europe and the United Kingdom. This makes South Africa's local resource one of the highest in the world." Citing http://www.energy.gov.za/files/esources/renewables/r_solar.html

Since our device needs to stay alert 24/7 for poles that might fall or short anytime, the power supply needs to be most reliable to the device. The device needs to be on for as long as possible. 

Another main concern was about the duality, and as long as there is sunshine there will be enough energy to supply power towards the system.  

\section{Additional Design Components}

\section{References}

\section{Appendices}

\end{document}